% SAX vs FR-SAX comparison figure (included via \input)
\usepackage{pgfplots}
\pgfplotsset{compat=1.18}
\usepackage{ctex}
\usetikzlibrary{patterns,decorations.pathreplacing,calligraphy}

\begin{document}

\begin{tikzpicture}

% ============================================================
% 生成一条模拟时间序列（归一化后），含一个瞬时尖峰
% 数据点（20个点，模拟归一化后的序列）
% ============================================================

% 定义数据
\def\data{
    0/0.3, 1/0.5, 2/0.8, 3/1.1, 4/0.9,
    5/0.6, 6/0.2, 7/-0.3, 8/-0.8, 9/-1.2,
    10/-0.9, 11/-0.4, 12/0.1, 13/0.7, 14/1.3,
    15/0.5, 16/-0.1, 17/-0.6, 18/-0.3, 19/0.2
}

% 分位数线 (alpha=5, 标准正态分位数)
\def\qA{-0.84}  % beta_1
\def\qB{-0.25}  % beta_2
\def\qC{0.25}   % beta_3
\def\qD{0.84}   % beta_4

% ============================================================
% 左图：传统SAX（含PAA分段聚合）
% ============================================================
\begin{scope}[xshift=0cm]

\begin{axis}[
    name=saxplot,
    width=8.5cm, height=6.5cm,
    xlabel={时间},
    ylabel={归一化值},
    title={\textbf{(a) 传统SAX算法}},
    title style={font=\small},
    xmin=-0.5, xmax=19.5,
    ymin=-1.6, ymax=1.6,
    xtick={0,4,8,12,16,19},
    ytick={-1.2,-0.84,-0.25,0.25,0.84,1.2},
    yticklabels={$-1.2$,$\beta_1$,$\beta_2$,$\beta_3$,$\beta_4$,$1.2$},
    grid=none,
    xlabel style={font=\small},
    ylabel style={font=\small},
    tick label style={font=\tiny},
    legend style={font=\tiny, at={(0.02,0.98)}, anchor=north west, draw=none, fill=white, fill opacity=0.8},
    clip=false,
]

% 分位数线
\draw[dashed, gray!60, thin] (axis cs:-0.5,-0.84) -- (axis cs:19.5,-0.84);
\draw[dashed, gray!60, thin] (axis cs:-0.5,-0.25) -- (axis cs:19.5,-0.25);
\draw[dashed, gray!60, thin] (axis cs:-0.5,0.25) -- (axis cs:19.5,0.25);
\draw[dashed, gray!60, thin] (axis cs:-0.5,0.84) -- (axis cs:19.5,0.84);

% 区域标注（右侧）
\node[font=\tiny, gray] at (axis cs:20.5, -1.22) {$l_1$};
\node[font=\tiny, gray] at (axis cs:20.5, -0.545) {$l_2$};
\node[font=\tiny, gray] at (axis cs:20.5, 0) {$l_3$};
\node[font=\tiny, gray] at (axis cs:20.5, 0.545) {$l_4$};
\node[font=\tiny, gray] at (axis cs:20.5, 1.22) {$l_5$};

% 原始曲线
\addplot[blue!60, thin, mark=*, mark size=1pt] coordinates {
    (0,0.3) (1,0.5) (2,0.8) (3,1.1) (4,0.9)
    (5,0.6) (6,0.2) (7,-0.3) (8,-0.8) (9,-1.2)
    (10,-0.9) (11,-0.4) (12,0.1) (13,0.7) (14,1.3)
    (15,0.5) (16,-0.1) (17,-0.6) (18,-0.3) (19,0.2)
};

% PAA分段（4段，每段5个点）
% 段1: (0.3+0.5+0.8+1.1+0.9)/5 = 0.72
% 段2: (0.6+0.2-0.3-0.8-1.2)/5 = -0.30
% 段3: (-0.9-0.4+0.1+0.7+1.3)/5 = 0.16
% 段4: (0.5-0.1-0.6-0.3+0.2)/5 = -0.06
\addplot[red!80, very thick, const plot] coordinates {
    (0,0.72) (5,-0.30) (10,0.16) (15,-0.06) (20,-0.06)
};

% 符号标注
\node[font=\small, red!80, above] at (axis cs:2.5, 0.72) {$l_4$};
\node[font=\small, red!80, below] at (axis cs:7.5, -0.30) {$l_2$};
\node[font=\small, red!80, above] at (axis cs:12.5, 0.16) {$l_3$};
\node[font=\small, red!80, below] at (axis cs:17.5, -0.06) {$l_3$};

% PAA段分隔线
\draw[gray, dashed, thin] (axis cs:5,-1.6) -- (axis cs:5,1.6);
\draw[gray, dashed, thin] (axis cs:10,-1.6) -- (axis cs:10,1.6);
\draw[gray, dashed, thin] (axis cs:15,-1.6) -- (axis cs:15,1.6);

\legend{归一化序列, PAA分段聚合}

% 底部符号序列
\node[font=\small, anchor=north] at (axis cs:9.5, -1.85) {符号序列：$l_4\ l_2\ l_3\ l_3$};

\end{axis}

\end{scope}

% ============================================================
% 右图：FR-SAX（无PAA，逐点映射）
% ============================================================
\begin{scope}[xshift=9cm]

\begin{axis}[
    name=frsaxplot,
    width=8.5cm, height=6.5cm,
    xlabel={时间},
    ylabel={归一化值},
    title={\textbf{(b) FR-SAX算法}},
    title style={font=\small},
    xmin=-0.5, xmax=19.5,
    ymin=-1.6, ymax=1.6,
    xtick={0,4,8,12,16,19},
    ytick={-1.2,-0.84,-0.25,0.25,0.84,1.2},
    yticklabels={$-1.2$,$\beta_1$,$\beta_2$,$\beta_3$,$\beta_4$,$1.2$},
    grid=none,
    xlabel style={font=\small},
    ylabel style={font=\small},
    tick label style={font=\tiny},
    legend style={font=\tiny, at={(0.02,0.98)}, anchor=north west, draw=none, fill=white, fill opacity=0.8},
    clip=false,
]

% 分位数线
\draw[dashed, gray!60, thin] (axis cs:-0.5,-0.84) -- (axis cs:19.5,-0.84);
\draw[dashed, gray!60, thin] (axis cs:-0.5,-0.25) -- (axis cs:19.5,-0.25);
\draw[dashed, gray!60, thin] (axis cs:-0.5,0.25) -- (axis cs:19.5,0.25);
\draw[dashed, gray!60, thin] (axis cs:-0.5,0.84) -- (axis cs:19.5,0.84);

% 区域标注（右侧）
\node[font=\tiny, gray] at (axis cs:20.5, -1.22) {$l_1$};
\node[font=\tiny, gray] at (axis cs:20.5, -0.545) {$l_2$};
\node[font=\tiny, gray] at (axis cs:20.5, 0) {$l_3$};
\node[font=\tiny, gray] at (axis cs:20.5, 0.545) {$l_4$};
\node[font=\tiny, gray] at (axis cs:20.5, 1.22) {$l_5$};

% 原始曲线
\addplot[blue!60, thin, mark=*, mark size=1pt] coordinates {
    (0,0.3) (1,0.5) (2,0.8) (3,1.1) (4,0.9)
    (5,0.6) (6,0.2) (7,-0.3) (8,-0.8) (9,-1.2)
    (10,-0.9) (11,-0.4) (12,0.1) (13,0.7) (14,1.3)
    (15,0.5) (16,-0.1) (17,-0.6) (18,-0.3) (19,0.2)
};

% FR-SAX逐点符号化（阶梯线，每个点映射到其所在区间的中心值）
% 0.3 -> l_4 (0.545)
% 0.5 -> l_4 (0.545)
% 0.8 -> l_4 (0.545)
% 1.1 -> l_5 (1.22)
% 0.9 -> l_5 (1.22)
% 0.6 -> l_4 (0.545)
% 0.2 -> l_3 (0)
% -0.3 -> l_2 (-0.545)
% -0.8 -> l_2 (-0.545)
% -1.2 -> l_1 (-1.22)
% -0.9 -> l_1 (-1.22)
% -0.4 -> l_2 (-0.545)
% 0.1 -> l_3 (0)
% 0.7 -> l_4 (0.545)
% 1.3 -> l_5 (1.22)
% 0.5 -> l_4 (0.545)
% -0.1 -> l_3 (0)
% -0.6 -> l_2 (-0.545)
% -0.3 -> l_2 (-0.545)
% 0.2 -> l_3 (0)

\addplot[green!60!black, very thick, const plot] coordinates {
    (0,0.545) (1,0.545) (2,0.545) (3,1.22) (4,1.22)
    (5,0.545) (6,0) (7,-0.545) (8,-0.545) (9,-1.22)
    (10,-1.22) (11,-0.545) (12,0) (13,0.545) (14,1.22)
    (15,0.545) (16,0) (17,-0.545) (18,-0.545) (19,0) (20,0)
};

\legend{归一化序列, 逐点符号映射}

% 底部符号序列
\node[font=\tiny, anchor=north, text width=8cm, align=center] at (axis cs:9.5, -1.85) {符号序列：$l_4\ l_4\ l_4\ l_5\ l_5\ l_4\ l_3\ l_2\ l_2\ l_1\ l_1\ l_2\ l_3\ l_4\ l_5\ l_4\ l_3\ l_2\ l_2\ l_3$};

\end{axis}

\end{scope}

\end{tikzpicture}

\end{document}
